\section{Architecture}
\label{sec:architecture}

This section presents the six biological principles underlying MI and the data structures that implement them. Each subsection maps a biological concept to a precise computational mechanism. Figure~\ref{fig:architecture} shows the overall dual-clock loop and how Endoquilibrium modulates each subsystem.

\begin{figure}[h]
  \centering
  \figifexists[\linewidth]{architecture.pdf}{architecture diagram}
  \caption{Morphon-Core dual-clock developmental loop. Solid arrows are data flow; dashed red arrows are Endoquilibrium modulation; dotted arrow is the structural-feedback path where slow-clock synaptogenesis introduces new connections that the fast clock then uses. The four temporal scales (fast/medium/slow/glacial) operate on the same morphon population at different periods.}
  \label{fig:architecture}
\end{figure}

\begin{figure}[h]
  \centering
  \figifexists[0.85\linewidth]{topology_developed.jpg}{Morphon network topology}
  \caption{The Morphon network during runtime, rendered by the in-browser WASM visualizer. Spheres are morphons colored by cell type: purple = associative, cyan = sensory, orange = motor, green = modulatory. Dotted yellow trails are spike-carrying synapses; the discrete dots represent action potentials propagating along axons with their learned delays. The 3D layout is the Poincar\'e ball embedding (Section~\ref{sec:architecture}); morphons of the same type cluster naturally because their hyperbolic positions are correlated with cell-type identity. Orange motor morphons concentrate on one side and project to the others through the synaptic web. The visualizer is reproducible from \texttt{web/} via \texttt{wasm-pack build} and a local HTTP server.}
  \label{fig:topology}
\end{figure}

\subsection{The Six Biological Principles}

\begin{description}
  \item[\textbf{P1: Local computation only.}] No global loss, no backpropagation. All updates use locally available information --- pre-synaptic and post-synaptic activity, locally diffused neuromodulators, and a Morphon's own state.
  \item[\textbf{P2: Developmental lifecycle.}] Morphons are born, differentiate, mature, fuse with correlated neighbors, and can die from energy starvation or chronic inactivity (apoptosis).
  \item[\textbf{P3: Chemical signaling.}] Morphogen-like signals diffuse through the embedding space and guide differentiation and connectivity decisions.
  \item[\textbf{P4: Neuromodulatory gating.}] Plasticity rules are gated by four modulatory channels (dopamine, serotonin, acetylcholine, norepinephrine analogues), each with a per-Morphon receptor density.
  \item[\textbf{P5: Multi-scale memory.}] Synaptic (fast, milliseconds), structural (medium, minutes), and morphogenetic (slow, hours) memory operate on separate timescales.
  \item[\textbf{P6: Metabolic cost.}] Every Morphon consumes energy each tick. Resource scarcity drives pruning, fusion, and apoptosis.
\end{description}

\subsection{The Morphon}

The Morphon is the fundamental compute unit. It carries:
\begin{itemize}
  \item \textbf{Membrane state:} potential, prediction-error baseline, refractory timer, fired flag.
  \item \textbf{Receptor profile:} per-modulator sensitivity vector $(\rho_R, \rho_N, \rho_A, \rho_H)$ for reward, novelty, arousal, and homeostasis channels.
  \item \textbf{Developmental metadata:} cell type (Stem, Sensory, Associative, Motor, Modulatory, Inhibitory Interneuron), differentiation level, lineage ID, birth tick.
  \item \textbf{Energy budget:} current energy, basal regeneration rate, firing cost (with optional superlinear penalty), apoptosis threshold.
  \item \textbf{Position:} hyperbolic coordinate in a Poincaré ball of configurable dimensionality.
\end{itemize}

A Morphon's state update is local: it integrates spike inputs from incoming synapses, applies an activation function determined by its cell type, fires if the result exceeds its (adaptive) threshold, and pays the metabolic cost.

\subsection{Resonance and Spike Propagation}

Morphons communicate via discrete spikes propagated along synapses with configurable delays. Unlike Transformer attention (a global $O(N^2)$ operation), spike propagation is local: $O(k \cdot N)$ where $k$ is the average out-degree. The system's cost scales with the number of \emph{active} morphons and their connectivity, not with the square of the population.

Each synapse carries a weight, a delay, an eligibility trace, a tag (for synaptic tagging and capture \cite{frey1997synaptic}), and a consolidation level. The dual-clock scheduler updates synaptic weights on the medium path via a three-factor rule:
\begin{equation}
\Delta w = \eta \cdot e \cdot M(t)
\end{equation}
where $e$ is the eligibility trace built from STDP coincidences, $M(t)$ is the (receptor-gated) modulatory signal, and $\eta$ is a per-synapse plasticity rate.

\subsection{Morphogenesis Lifecycle}

The morphogenesis subsystem implements Principles P2 and P3. On the slow clock:
\begin{itemize}
  \item \textbf{Synaptogenesis} forms new synapses between morphons whose activity is correlated and whose hyperbolic distance is below threshold.
  \item \textbf{Pruning} removes synapses with low weight, low recent activity, or whose pre/post morphons are no longer correlated.
  \item \textbf{Migration} moves morphons through hyperbolic space along gradients of correlated input (similar morphons drift together).
\end{itemize}
On the glacial clock:
\begin{itemize}
  \item \textbf{Division} produces a daughter morphon from a parent that is consistently energetic and active.
  \item \textbf{Differentiation} commits a stem morphon to a specific cell type based on local morphogen concentrations.
  \item \textbf{Fusion} merges two morphons whose activity patterns have become indistinguishable, redistributing their synapses to the survivor.
  \item \textbf{Apoptosis} removes morphons that have been below their energy threshold for a configurable duration.
\end{itemize}

\subsection{Endoquilibrium: Predictive Neuroendocrine Regulation}

Endoquilibrium is the homeostatic regulator. It monitors seven vital signs (firing rates by cell type, eligibility density, weight entropy, energy utilization, prediction error mean, reward EMA), maintains dual-timescale exponential moving averages (fast $\tau=50$, slow $\tau=500$), and adjusts six modulation channels through proportional control: $\rho_R, \rho_N, \rho_A, \rho_H$ gains, threshold bias, plasticity multiplier, and consolidation gain.

Endoquilibrium also detects developmental stage: \emph{Proliferating}, \emph{Differentiating}, \emph{Consolidating}, \emph{Mature}, \emph{Stressed}. The stage detection logic uses reward-based relative thresholds and is the locus of the regulatory failure mode we document in Section~\ref{sec:failure_modes}.

\subsection{Hyperbolic Embedding Space}

Morphons live in a Poincaré ball model of hyperbolic space with configurable dimensionality (typically 4--6). Hyperbolic geometry has two properties that make it suitable for developmental architectures: (i) tree-like hierarchies are represented with low distortion \cite{nickel2017poincare}, and (ii) volume grows exponentially with radius, providing room for unbounded topology expansion without coordinate crowding.

In MI, the origin of the ball corresponds to the most general (stem-like) morphons, and the boundary to the most specialized. A morphon's hyperbolic distance to its peers governs synaptogenesis probability and morphogen diffusion.

\subsection{Dual-Clock Scheduler}

Four temporal scales separate fast inference from slow structural change:
\begin{center}
\begin{tabular}{lll}
\toprule
Scale & Period & Operations \\
\midrule
Fast    & 1 step    & Spike propagation, firing, input integration \\
Medium  & 5--10     & STDP, eligibility, three-factor weight updates \\
Slow    & 100--1000 & Synaptogenesis, pruning, migration \\
Glacial & 1000--5000 & Division, differentiation, fusion, apoptosis \\
\bottomrule
\end{tabular}
\end{center}
This separation is essential. Structural changes during inference would destabilize the network; the slow clock guarantees that morphogenetic events happen between bursts of computation.
